% !TeX root = ../../main.tex
% =====================================================================
%  Fundamentos de mecánica aplicada
%  Responsable:
%  Última revisión:
% =====================================================================
\chapter{Fundamentos de mecánica aplicada}
\label{ch:mecanica_aplicada}

\begin{objetivos}
  \item Plantear el diagrama de sólido libre de una esquina del coche y sacar de
        él las cargas de cada barra.
  \item Identificar si una pieza está limitada por resistencia, por rigidez o por
        pandeo, y dimensionarla en consecuencia.
  \item Saber cuándo hay que comprobar fatiga y cuándo no merece la pena.
  \item Calcular el par de apriete de un tornillo a partir de la precarga que se
        quiere conseguir, y entender por qué la precarga es lo único que importa.
  \item Elegir material con criterio, sabiendo qué diferencias entre materiales
        son reales y cuáles son mitología.
\end{objetivos}

\fichasesion{90 minutos}{Resistencia de materiales de primer ciclo}%
{Calculadora, hoja de cálculo, un tornillo y una llave dinamométrica}

Este capítulo no sustituye a una asignatura de resistencia de materiales: da por
sabido lo básico y se centra en las cinco cosas que un equipo de Formula Student
usa constantemente y suele hacer mal. Para profundizar,
\textcite{budynas_shigley_2020} cubre todo lo que hay aquí con mucho más detalle.

% ---------------------------------------------------------------------
\section{Diagramas de sólido libre}
\label{sec:mecanica_aplicada:01}
\index{diagrama de sólido libre}

Todo cálculo estructural empieza aislando un trozo del coche y dibujando todo lo
que actúa sobre él. Parece elemental, pero conviene hacerlo con cuidado: un
diagrama mal planteado da un resultado coherente consigo mismo, y por eso el
error no se detecta después.

El método no tiene misterio:

\begin{enumerate}
  \item \textbf{Aísla} el cuerpo y dibújalo separado de todo lo demás.
  \item \textbf{Sustituye} cada cosa que has quitado por las fuerzas y momentos
        que ejercía.
  \item \textbf{Aplica} equilibrio: $\sum \vect{F} = \vect{0}$ y
        $\sum \vect{M} = \vect{0}$ respecto de un punto cualquiera.
  \item \textbf{Comprueba}: la suma de reacciones debe dar la carga aplicada.
\end{enumerate}

\subsection{El caso canónico: una esquina del coche}

El ejercicio que hay que saber hacer de memoria es este: dada una fuerza en la
huella del neumático, obtener la carga axial de cada barra de la suspensión.

La simplificación que lo hace tratable es que \textbf{los trapecios, el tirante
de dirección y el push\-rod son barras biarticuladas}: rotuladas en ambos
extremos, no pueden transmitir momento, así que solo trabajan a tracción o
compresión pura. Por eso son tubos y por eso el cálculo se reduce a resolver un
sistema de equilibrio de fuerzas concurrentes.

La fuerza de entrada actúa en la \emph{huella de contacto}, no en el centro de
rueda. La diferencia es el radio del neumático multiplicado por la fuerza
lateral, y es un momento que hay que llevar hasta la mangueta.

\begin{ojo}
\begin{itemize}
  \item Aplicar la carga en el centro de rueda en vez de en la huella.
  \item Olvidar el equilibrio de momentos: con solo fuerzas, el sistema parece
        resuelto y no lo está.
  \item Suponer que una barra con casquillos rígidos (no rotulados) trabaja solo
        a axial. Si el extremo no gira libremente, hay flexión.
  \item No comprobar el resultado: si las reacciones no suman la carga aplicada,
        hay un error. Esta comprobación cuesta treinta segundos.
\end{itemize}
\end{ojo}

% ---------------------------------------------------------------------
\section{Tensión, rigidez y pandeo}
\label{sec:mecanica_aplicada:02}
\index{pandeo}

\subsection{Las tres preguntas}

Antes de dimensionar cualquier pieza hay que saber qué la limita, porque cada
límite se combate de forma distinta y a veces con soluciones opuestas:

\begin{center}
\begin{tabularx}{\textwidth}{@{}lXl@{}}
\toprule
\textbf{Limitada por} & \textbf{Pregunta} & \textbf{Qué la mejora} \\
\midrule
Resistencia & ¿Rompe o plastifica? & Material más resistente, más sección \\
Rigidez     & ¿Se deforma demasiado? & Más momento de inercia, mejor forma \\
Pandeo      & ¿Colapsa a compresión? & Más diámetro, menos longitud libre \\
Fatiga      & ¿Rompe tras muchos ciclos? & Radios, acabado, menos concentración \\
\bottomrule
\end{tabularx}
\end{center}

En un coche de Formula Student, \textbf{la mayoría de las piezas estructurales no
están limitadas por resistencia}. Los trapecios los limita el pandeo, el chasis
lo limita la rigidez torsional, y los elementos que giran los limita la fatiga.
Dimensionar todo contra el límite elástico produce un coche pesado que además
falla por donde no se miraba.

\subsection{Tensiones básicas}

\begin{align}
  \sigma_{\text{axial}} &= \frac{N}{A} &
  \sigma_{\text{flexión}} &= \frac{M\,y}{I} &
  \tau_{\text{torsión}}  &= \frac{T\,r}{J}
  \label{eq:tensiones_basicas}
\end{align}

Cuando actúan varias a la vez, se combinan con el criterio de von Mises, que es
el que se compara con el límite elástico:

\begin{equation}
  \sigma_{VM} = \sqrt{\sigma_x^2 - \sigma_x\sigma_y + \sigma_y^2 + 3\tau_{xy}^2}
  \label{eq:von_mises}
\end{equation}

\subsection{Pandeo: lo que gobierna los trapecios}

Una barra esbelta a compresión no falla por aplastamiento: colapsa lateralmente
mucho antes. La carga crítica de Euler es

\begin{equation}
  P_{cr} = \frac{\pi^2 E I}{(K L)^2}
  \label{eq:euler}
\end{equation}

donde $L$ es la longitud entre apoyos y $K$ depende de cómo estén sujetos los
extremos:

\begin{center}
\begin{tabular}{@{}lcl@{}}
\toprule
\textbf{Condición de extremos} & $K$ & \textbf{Caso típico} \\
\midrule
Articulado--articulado & 1{,}0 & \textbf{Barra de suspensión con rótulas} \\
Empotrado--articulado  & 0{,}7 & Extremo soldado, otro rotulado \\
Empotrado--empotrado   & 0{,}5 & Barra soldada en ambos extremos \\
Empotrado--libre       & 2{,}0 & Voladizo \\
\bottomrule
\end{tabular}
\end{center}

Para el caso normal en suspensión --- rótula en los dos extremos --- $K = 1$, que
es el peor caso razonable. Conviene usarlo aunque haya algo de empotramiento: la
diferencia entre $K=1$ y $K=0{,}7$ duplica la carga crítica, y esa no es una
apuesta que merezca la pena hacer.

\begin{clave}
Para un tubo de pared delgada de radio medio $r_m$ y espesor $t$:
$A \approx 2\pi r_m t$ e $I \approx \pi r_m^3 t$. Sustituyendo en
\eqref{eq:euler} con la \emph{masa constante} (es decir, $r_m t$ fijo), resulta
$P_{cr} \propto r_m^2$.

Traducido: \textbf{duplicar el diámetro de un tubo, adelgazando la pared para
mantener el peso, multiplica por cuatro la carga de pandeo.} Por eso los
trapecios son tubos de diámetro grande y pared fina, y por eso el límite práctico
no lo pone el cálculo sino la pared mínima que se puede soldar, el pandeo local
de la pared y lo que exija el reglamento.
\end{clave}

La fórmula de Euler solo vale para barras esbeltas. Para esbelteces bajas
($\lambda = KL/r_g$ pequeño, con $r_g=\sqrt{I/A}$ el radio de giro) el fallo pasa
a ser plastificación y hay que usar una curva de transición --- por ejemplo la
parábola de Johnson. En la práctica: si la carga de Euler sale por encima de la
carga de plastificación $A\,S_y$, el pandeo no es el modo crítico y hay que
comprobar resistencia.

% ---------------------------------------------------------------------
\section{Fatiga}
\label{sec:mecanica_aplicada:03}
\index{fatiga}

Una pieza sometida a carga cíclica puede romper con tensiones muy por debajo del
límite elástico. La grieta nace en la superficie, casi siempre en un cambio de
sección, y crece hasta que la sección restante no aguanta.

\subsection{¿Importa la fatiga en un coche de Formula Student?}

Depende de la pieza, y merece pensarlo antes de gastar tiempo:

\begin{center}
\begin{tabularx}{\textwidth}{@{}lXl@{}}
\toprule
\textbf{Tipo de pieza} & \textbf{Por qué} & \textbf{¿Fatiga?} \\
\midrule
Palieres, bujes, tornillos de rueda & Giran: miles de ciclos por kilómetro & \textbf{Sí, siempre} \\
Anclajes del chasis, trapecios & Carga variable, muchos km de test & Sí \\
Uniones soldadas & El cordón concentra tensión & Sí \\
Piezas que se reutilizan otra temporada & Acumulan historia desconocida & Sí \\
Estructura cargada casi estáticamente & Pocos ciclos de amplitud alta & Normalmente no \\
\bottomrule
\end{tabularx}
\end{center}

La resistencia de \SI{22}{\kilo\meter} induce a pensar que la fatiga no importa.
Es engañoso por tres motivos: la temporada de ensayos son cientos de kilómetros,
las piezas que giran ven un ciclo completamente invertido por vuelta de rueda, y
las piezas heredadas de temporadas anteriores llegan con un daño acumulado que
nadie ha contabilizado.

\subsection{Límite de fatiga}

Para aceros, el límite de fatiga en probeta se estima como
$S'_e \approx 0{,}5\,S_{ut}$ (con un techo en torno a \SI{700}{\mega\pascal}). Ese
valor de laboratorio se corrige con los factores de Marin para llegar al de la
pieza real:

\begin{equation}
  S_e = k_a\,k_b\,k_c\,k_d\,k_e\; S'_e
  \label{eq:marin}
\end{equation}

donde $k_a$ recoge el acabado superficial, $k_b$ el tamaño, $k_c$ el tipo de
carga, $k_d$ la temperatura y $k_e$ la fiabilidad exigida. El más castigador
suele ser $k_a$: una superficie bruta de laminación puede reducir el límite de
fatiga a menos de la mitad respecto de una pulida.

\begin{ojo}
El aluminio \textbf{no tiene límite de fatiga}: su curva S--N sigue bajando
indefinidamente. Una pieza de aluminio no se puede diseñar «para vida infinita»;
hay que fijar un número de ciclos objetivo y una vida en kilómetros, y retirarla
o inspeccionarla al alcanzarla. Esto afecta directamente a manguetas, balancines
y anclajes mecanizados.
\end{ojo}

\subsection{Concentración de tensiones y tensión media}

El efecto de una entalla se recoge con el factor de concentración $K_t$
(geométrico) corregido por la sensibilidad a la entalla $q$ del material:

\begin{equation}
  K_f = 1 + q\,(K_t - 1)
  \label{eq:kf}
\end{equation}

Y cuando la carga cíclica se superpone a una carga media --- lo habitual --- se
comprueba con un criterio como el de Goodman:

\begin{equation}
  \frac{\sigma_a}{S_e} + \frac{\sigma_m}{S_{ut}} = \frac{1}{n}
  \label{eq:goodman}
\end{equation}

\begin{clave}
En fatiga, la geometría manda sobre el material. Un radio de acuerdo generoso, un
buen acabado superficial y evitar cambios bruscos de sección hacen más por la
vida de una pieza que subir de grado de material --- que además, al ser más
resistente, suele ser \emph{más} sensible a la entalla.
\end{clave}

% ---------------------------------------------------------------------
\section{Uniones atornilladas y pares de apriete}
\label{sec:mecanica_aplicada:04}
\index{uniones atornilladas}\index{par de apriete}

Buena parte de los fallos de un coche de competición aparecen en las uniones
antes que en las piezas, y muchos de ellos se explican por una precarga
incorrecta.

\subsection{La precarga lo es todo}

Un tornillo bien apretado no es un tornillo que sujeta: es un tornillo que
\emph{comprime} las piezas entre sí con una fuerza mucho mayor que la carga
externa. Mientras esa compresión no desaparezca, la unión no se mueve, no vibra y
--- esto es lo importante --- el tornillo apenas ve la carga variable.

La precarga recomendada es una fracción alta de la carga de prueba del tornillo:

\begin{equation}
  F_i = 0{,}75\, A_t\, S_p \quad\text{(uniones desmontables)}, \qquad
  F_i = 0{,}90\, A_t\, S_p \quad\text{(uniones permanentes)}
  \label{eq:precarga}
\end{equation}

con $A_t$ el área resistente a tracción y $S_p$ la carga de prueba del grado:

\begin{center}
\begin{tabular}{@{}lSS@{}}
\toprule
\textbf{Grado} & {$S_p$ (\si{\mega\pascal})} & {$S_{ut}$ (\si{\mega\pascal})} \\
\midrule
8.8  & 600 & 800 \\
10.9 & 830 & 1040 \\
12.9 & 970 & 1220 \\
\bottomrule
\end{tabular}
\end{center}

\subsection{Del par de apriete a la precarga}

La relación práctica es

\begin{equation}
  T = K\, d\, F_i
  \label{eq:par_apriete}
\end{equation}

con $d$ el diámetro nominal y $K$ un coeficiente que vale aproximadamente
\num{0.20} en seco y \num{0.15} lubricado.

\begin{ojo}
El coeficiente $K$ tiene una dispersión enorme --- fácilmente $\pm\SI{25}{\percent}$
--- porque depende del rozamiento en la rosca y bajo la cabeza, que a su vez
depende del acabado, la lubricación y si el tornillo es nuevo. Es decir:
\textbf{una llave dinamométrica controla el par, no la precarga}. Es lo mejor que
tenemos en el taller, pero hay que saber que el número real de precarga tiene esa
incertidumbre encima, y no diseñar uniones que dependan de acertarla con precisión.
\end{ojo}

\subsection{Reparto de la carga externa}

La rigidez relativa del tornillo $k_b$ y de las piezas $k_m$ determina qué parte
de una carga externa $P$ ve el tornillo:

\begin{equation}
  C = \frac{k_b}{k_b + k_m}, \qquad
  F_{\text{tornillo}} = F_i + C\,P, \qquad
  F_{\text{unión}} = F_i - (1-C)\,P
  \label{eq:reparto_carga}
\end{equation}

Como las piezas suelen ser mucho más rígidas que el tornillo, $C$ es pequeño
(típicamente \numrange{0.1}{0.3}) y el tornillo apenas nota la carga variable. De
ahí una regla de diseño poco intuitiva: \textbf{un tornillo largo (menos rígido)
aguanta mejor a fatiga que uno corto}, porque reduce $C$.

La unión se separa cuando $F_i - (1-C)P = 0$. Diseñar con margen frente a ese
punto es más importante que comprobar la tensión del tornillo.

\subsection{Reglas prácticas}

\begin{itemize}
  \item \textbf{Cortante: por aplastamiento, no por rozamiento.} Salvo que la
        precarga esté controlada de verdad, se dimensiona suponiendo que el
        tornillo trabaja a cortante contra el agujero. Y la rosca nunca debe
        quedar en el plano de cortadura: se usa tornillo de vástago liso.
  \item \textbf{Longitud de rosca engranada}: al menos $1\,d$ en acero y
        $2\,d$ en aluminio. En aluminio, mejor inserto roscado.
  \item \textbf{Arandelas} siempre bajo la cabeza que gira, y bajo cualquier
        apoyo en aluminio.
  \item \textbf{No reutilizar} tornillos apretados cerca del límite elástico, ni
        tuercas autoblocantes de anillo plástico que ya hayan trabajado.
  \item \textbf{Marcar el apriete}: una raya de pintura entre tornillo y pieza
        permite ver de un vistazo si algo se ha aflojado, tanto en el taller como
        en la revisión técnica.
\end{itemize}

\begin{regla}[frenado de tornillería]
El reglamento exige \emph{frenado positivo} en los tornillos críticos: elementos
que impidan mecánicamente el aflojamiento, como tuercas autoblocantes, tuerca
almenada con pasador o alambre de frenado. Los adhesivos de bloqueo de roscas
habitualmente \textbf{no} se aceptan por sí solos como frenado positivo.
Compruébalo en el artículo correspondiente del reglamento vigente antes de
diseñar la unión: es un motivo clásico de rechazo en la revisión técnica.
\end{regla}

% ---------------------------------------------------------------------
\section{Selección de materiales}
\label{sec:mecanica_aplicada:05}
\index{selección de materiales}

\subsection{El dato que casi nadie espera}

La rigidez específica --- módulo elástico dividido por densidad --- de casi todos
los metales estructurales es prácticamente la misma:

\begin{center}
\begin{tabular}{@{}lSSS@{}}
\toprule
\textbf{Material} & {$E$ (\si{\giga\pascal})} & {$\rho$ (\si{\gram\per\cubic\centi\meter})} & {$E/\rho$} \\
\midrule
Acero (cualquier aleación) & 210 & 7.85 & 26.8 \\
Aluminio                   & 70  & 2.70 & 25.9 \\
Titanio                    & 110 & 4.51 & 24.4 \\
Magnesio                   & 45  & 1.74 & 25.9 \\
\bottomrule
\end{tabular}
\end{center}

\begin{clave}
Dos consecuencias que ahorran muchas discusiones:

\begin{enumerate}
  \item \textbf{Un acero aleado no es más rígido que un acero normal.} Todos los
        aceros tienen $E \approx \SI{210}{\giga\pascal}$. Se elige un acero
        aleado por resistencia, no por rigidez.
  \item \textbf{A igualdad de masa, y si la pieza trabaja a tracción pura, da
        igual el metal.} La ventaja del aluminio aparece cuando la pieza trabaja a
        \emph{flexión o pandeo}: como es menos denso, con la misma masa se hace
        más voluminosa, y el momento de inercia crece mucho más deprisa que la
        pérdida de módulo. Es un efecto de \emph{forma}, no de material.
\end{enumerate}
\end{clave}

Donde sí hay diferencias enormes es en la resistencia específica, y ahí es donde
la elección de aleación y tratamiento importa de verdad.

\subsection{Paleta habitual en Formula Student}

\begin{center}
\begin{tabularx}{\textwidth}{@{}llX@{}}
\toprule
\textbf{Familia} & \textbf{Designación} & \textbf{Cuándo y avisos} \\
\midrule
Acero al carbono & S235, S355 & Utillaje, soportes poco cargados. Barato y disponible \\
Acero aleado & 25CrMo4 & Tubos de chasis y piezas soldadas muy cargadas. Soldable \\
Acero aleado & 42CrMo4 & Ejes y palieres. Requiere tratamiento térmico \\
Aluminio & 6082-T6 & Uso general mecanizado. \textbf{Soldable}, buena corrosión \\
Aluminio & 7075-T6 & Manguetas, balancines. Alta resistencia. \textbf{No soldable}, mala corrosión \\
Titanio & Ti-6Al-4V & Solo si está justificado con números. Caro y difícil de mecanizar \\
\bottomrule
\end{tabularx}
\end{center}

\pendiente{Ajustar esta tabla a lo que el equipo puede comprar y mecanizar de
verdad, con proveedores y plazos. Un material que no se puede conseguir en dos
semanas en el formato necesario no es una opción de diseño.}

\subsection{Método de selección}

\begin{enumerate}
  \item \textbf{Identifica el modo limitante} (resistencia, rigidez, pandeo,
        fatiga) y la geometría disponible. Este paso decide casi todo.
  \item \textbf{Escribe el índice de material} que corresponde a ese modo
        (\textcite{ashby_materials_2016} los tabula todos) y ordena candidatos.
  \item \textbf{Filtra por fabricabilidad}: ¿se puede soldar, mecanizar, tratar
        con lo que tenemos o con lo que nos pueden hacer?
  \item \textbf{Filtra por disponibilidad y coste}: formatos comerciales, plazo de
        entrega, y lo que costará en el informe del capítulo
        \ref{ch:evento_costes}.
  \item \textbf{Comprueba compatibilidad}: contacto galvánico entre metales
        distintos, temperatura de servicio, y compatibilidad con adhesivos o
        fluidos (capítulo \ref{ch:fabricacion_ii}).
\end{enumerate}

\begin{ojo}
\begin{itemize}
  \item Elegir el material antes que la forma. El material es la última decisión,
        no la primera.
  \item Usar 7075 y descubrir en el taller que hay que soldarlo.
  \item Especificar un tratamiento térmico y no comprobar quién lo va a hacer,
        cuánto tarda y cuánto deforma la pieza.
  \item Poner aluminio y acero en contacto directo en una zona que se moja.
  \item Copiar el material de la pieza equivalente de otro equipo sin saber qué
        modo de fallo estaban combatiendo.
\end{itemize}
\end{ojo}

% ---------------------------------------------------------------------
%  Cierre del capítulo
% ---------------------------------------------------------------------

\begin{caso}
\redactar{Rellenar con nuestros casos reales: qué se ha roto en el coche y por
qué. Para cada fallo, indicar el modo (estático, pandeo, fatiga, unión floja), si
se había calculado y con qué caso de carga. Es el recuadro más valioso del
capítulo y solo se puede escribir mirando piezas rotas.}
\end{caso}

\begin{ojo}
\begin{itemize}
  \item Dimensionar todo contra el límite elástico sin preguntarse qué limita
        realmente la pieza.
  \item Aplicar un coeficiente de seguridad genérico a un caso de carga que nadie
        ha definido: el coeficiente no arregla que la carga esté mal.
  \item Apretar «a sensación». La diferencia entre un tornillo M8 bien apretado y
        uno apretado con fuerza es de más del doble de precarga.
  \item No documentar los pares de apriete: si no están en una tabla, en el
        montaje cada uno usa el suyo.
\end{itemize}
\end{ojo}

\begin{entregable}
\begin{enumerate}
  \item Una hoja de cálculo de \textbf{pandeo de tubos}: entradas diámetro,
        espesor, longitud y condiciones de extremo; salida carga crítica y
        coeficiente de seguridad.
  \item Una hoja de cálculo de \textbf{uniones atornilladas}: entradas grado,
        métrica y tipo de unión; salidas precarga, par de apriete y carga de
        separación.
  \item La \textbf{tabla de pares de apriete del coche}, por tipo de unión,
        impresa y colgada en el taller.
\end{enumerate}
\end{entregable}

\begin{juez}
\begin{enumerate}
  \item ¿Qué limita el dimensionado de este trapecio: resistencia, rigidez o pandeo?
  \item ¿Por qué este tubo tiene este diámetro y este espesor y no otros?
  \item ¿Con qué par se aprieta esto y de dónde sale ese número?
  \item ¿Has comprobado fatiga en alguna pieza? ¿En cuáles y por qué en esas?
  \item ¿Por qué esta pieza es de 7075 y no de 6082?
\end{enumerate}
\end{juez}
